% GENERATED FILE. Edit canonical lessons, then run npm run generate:books.
% canonical-source-hash: fnv1a64:95786b6e2660e2a5
% canonical-lessons: RU-C391-zakonodatelstvo, RU-C391-konstituciya, RU-C391-grazhdanstvo, RU-C391-yuridicheskiy, RU-C391-ogranichenie

\chapter{Laws of the land}
\label{ch:ru-laws-of-the-land}
\hlchaptermodality{391}

\begin{chapteropening}
\textbf{By the end of this chapter:} \emph{I can talk about legislation, a constitution and citizenship, legal help, and a restriction.}

Complete the last lesson of chapter 391: I can talk about legislation, a constitution and citizenship, legal help, and a restriction.
\end{chapteropening}

\section[\texorpdfstring{\ru{законодательство} (zakonodátelstvo)}{законодательство (zakonodátelstvo)}]{\ru{законодательство} (zakonodátelstvo) — legislation, the laws (of a country)}

\label{lesson:RU-C391-zakonodatelstvo}



\begin{quote}
Before the new one: say the Russian for a right, then the Russian for responsibility.
\end{quote}

\subsection*{You'll want to know: \ru{законодательство}}
\textbf{\ru{законодательство}} (\emph{zakonodátelstvo}) — \textquotedblleft{}legislation, the laws (of a country)\textquotedblright{}.

\begin{quote}
\ru{В} \ru{законодательстве} \ru{есть} \ru{новые} \ru{правила}.
\end{quote}

\begin{quote}
\emph{V zakonodátelstve yest nóvyye právila.}
\end{quote}

\begin{quote}
There are new rules in the legislation.
\end{quote}

Neuter. All the laws of a country taken together.

\begin{cousinweb}
\textbf{\ru{законодательство}} is built from \emph{\ru{закон}} \textquotedblleft{}a law\textquotedblright{} and \emph{-\ru{датель}}, from \emph{\ru{дать}} \textquotedblleft{}to give\textquotedblright{}, with the suffix \emph{-\ru{ство}}: the giving of laws.
\end{cousinweb}

\begin{grammarlens}[title={What you've built}]
One more word for telling, explaining and describing at B1.
\end{grammarlens}

\subsection*{Your turn}
\begin{itemize}
\raggedright
  \item \emph{Say it:} \emph{zakonodátelstvo}
  \item \emph{Say it:} \emph{zakonodátelstvo}, once more
  \item \emph{Say it:} a sentence of your own with \emph{zakonodátelstvo}
  \item \emph{Recall:} say the Russian for a right, then the Russian for responsibility, then say \emph{zakonodátelstvo} again
\end{itemize}

\begin{tcolorbox}[breakable,colback=teal!4,colframe=teal!35!black,title={Before you move on}]
What does \ru{законодательство} mean? (\textquotedblleft{}Legislation, the laws (of a country)\textquotedblright{}.) Say it once more.
\end{tcolorbox}
\section[\texorpdfstring{\ru{конституция} (konstitútsiya)}{конституция (konstitútsiya)}]{\ru{конституция} (konstitútsiya) — a constitution}

\label{lesson:RU-C391-konstituciya}



\begin{quote}
Before the new one: say the Russian for responsibility, then the Russian for legislation.
\end{quote}

\subsection*{You'll want to know: \ru{конституция}}
\textbf{\ru{конституция}} (\emph{konstitútsiya}) — \textquotedblleft{}a constitution\textquotedblright{}.

\begin{quote}
\ru{Конституция} — \ru{главный} \ru{закон} \ru{страны}.
\end{quote}

\begin{quote}
\emph{Konstitútsiya — glávnyy zakón straný.}
\end{quote}

\begin{quote}
The constitution is the main law of the country.
\end{quote}

Feminine. It also means a person’s build.

\begin{cousinweb}
\textbf{\ru{конституция}} is usually traced, through western languages, to Latin \emph{cōnstitūtiō} \textquotedblleft{}an arrangement, a decree\textquotedblright{} (Lewis \& Short), from \emph{cōnstituere} \textquotedblleft{}to set up, to establish\textquotedblright{}.
\end{cousinweb}

\begin{grammarlens}[title={What you've built}]
One more word for telling, explaining and describing at B1.
\end{grammarlens}

\subsection*{Your turn}
\begin{itemize}
\raggedright
  \item \emph{Say it:} \emph{konstitútsiya}
  \item \emph{Say it:} \emph{konstitútsiya}, once more
  \item \emph{Say it:} a sentence of your own with \emph{konstitútsiya}
  \item \emph{Recall:} say the Russian for responsibility, then the Russian for legislation, then say \emph{konstitútsiya} again
\end{itemize}

\begin{tcolorbox}[breakable,colback=teal!4,colframe=teal!35!black,title={Before you move on}]
What does \ru{конституция} mean? (\textquotedblleft{}A constitution\textquotedblright{}.) Say it once more.
\end{tcolorbox}
\section[\texorpdfstring{\ru{гражданство} (grazhdánstvo)}{гражданство (grazhdánstvo)}]{\ru{гражданство} (grazhdánstvo) — citizenship}

\label{lesson:RU-C391-grazhdanstvo}



\begin{quote}
Before the new one: say the Russian for legislation, then the Russian for a constitution.
\end{quote}

\subsection*{You'll want to know: \ru{гражданство}}
\textbf{\ru{гражданство}} (\emph{grazhdánstvo}) — \textquotedblleft{}citizenship\textquotedblright{}.

\begin{quote}
\ru{У} \ru{неё} \ru{двойное} \ru{гражданство}.
\end{quote}

\begin{quote}
\emph{U neyó dvoynóye grazhdánstvo.}
\end{quote}

\begin{quote}
She has dual citizenship.
\end{quote}

A neuter noun. \emph{\ru{Двойное} \ru{гражданство}} — dual citizenship. From \emph{\ru{гражданин}}, a citizen.

\begin{cousinweb}
\textbf{\ru{гражданство}} is built on \emph{\ru{гражданин}} \textquotedblleft{}a citizen\textquotedblright{}, on \emph{\ru{град}}, the Church Slavonic form of native \emph{\ru{город}} \textquotedblleft{}a town\textquotedblright{}, Proto-Slavic \emph{*gordŭ}. The \emph{-\ru{ра}-} is the Church Slavonic shape, \emph{-\ru{оро}-} the native one.
\end{cousinweb}

\begin{grammarlens}[title={What you've built}]
One more word for telling, explaining and describing at B1.
\end{grammarlens}

\subsection*{Your turn}
\begin{itemize}
\raggedright
  \item \emph{Say it:} \emph{grazhdánstvo}
  \item \emph{Say it:} \emph{grazhdánstvo}, once more
  \item \emph{Say it:} a sentence of your own with \emph{grazhdánstvo}
  \item \emph{Recall:} say the Russian for legislation, then the Russian for a constitution, then say \emph{grazhdánstvo} again
\end{itemize}

\begin{tcolorbox}[breakable,colback=teal!4,colframe=teal!35!black,title={Before you move on}]
What does \ru{гражданство} mean? (\textquotedblleft{}Citizenship\textquotedblright{}.) Say it once more.
\end{tcolorbox}
\section[\texorpdfstring{\ru{юридический} (yuridícheskiy)}{юридический (yuridícheskiy)}]{\ru{юридический} (yuridícheskiy) — of the law, legal}

\label{lesson:RU-C391-yuridicheskiy}



\begin{quote}
Before the new one: say the Russian for a constitution, then the Russian for citizenship.
\end{quote}

\subsection*{You'll want to know: \ru{юридический}}
\textbf{\ru{юридический}} (\emph{yuridícheskiy}) — \textquotedblleft{}of the law, legal\textquotedblright{}.

\begin{quote}
\ru{Мне} \ru{нужна} \ru{юридическая} \ru{помощь}.
\end{quote}

\begin{quote}
\emph{Mne nuzhná yuridícheskaya pómoshch.}
\end{quote}

\begin{quote}
I need legal help.
\end{quote}

\emph{\ru{Юридический} \ru{факультет}} — a law faculty.

\begin{cousinweb}
\textbf{\ru{юридический}} is usually traced to Latin \emph{iūridicus} \textquotedblleft{}of the administering of justice\textquotedblright{}, from \emph{iūs} \textquotedblleft{}right, law\textquotedblright{} and \emph{dīcere} \textquotedblleft{}to say\textquotedblright{} (Lewis \& Short).
\end{cousinweb}

\begin{grammarlens}[title={What you've built}]
One more word for telling, explaining and describing at B1.
\end{grammarlens}

\subsection*{Your turn}
\begin{itemize}
\raggedright
  \item \emph{Say it:} \emph{yuridícheskiy}
  \item \emph{Say it:} \emph{yuridícheskiy}, once more
  \item \emph{Say it:} a sentence of your own with \emph{yuridícheskiy}
  \item \emph{Recall:} say the Russian for a constitution, then the Russian for citizenship, then say \emph{yuridícheskiy} again
\end{itemize}

\begin{tcolorbox}[breakable,colback=teal!4,colframe=teal!35!black,title={Before you move on}]
What does \ru{юридический} mean? (\textquotedblleft{}Of the law, legal\textquotedblright{}.) Say it once more.
\end{tcolorbox}
\section[\texorpdfstring{\ru{ограничение} (ogranichéniye)}{ограничение (ogranichéniye)}]{\ru{ограничение} (ogranichéniye) — a restriction, a limit}

\label{lesson:RU-C391-ogranichenie}



\begin{quote}
Before the new one: say the Russian for citizenship, then the Russian for of the law.
\end{quote}

\subsection*{You'll want to know: \ru{ограничение}}
\textbf{\ru{ограничение}} (\emph{ogranichéniye}) — \textquotedblleft{}a restriction, a limit\textquotedblright{}.

\begin{quote}
\ru{Есть} \ru{ограничения} \ru{по} \ru{возрасту}.
\end{quote}

\begin{quote}
\emph{Yest ogranichéniya po vózrastu.}
\end{quote}

\begin{quote}
There are age restrictions.
\end{quote}

Neuter. \emph{\ru{Ограничение} \ru{скорости}} — a speed limit.

\begin{cousinweb}
\textbf{\ru{ограничение}} is built on \emph{\ru{ограничить}} \textquotedblleft{}to restrict\textquotedblright{}, from \emph{\ru{граница}} \textquotedblleft{}a border\textquotedblright{}, which is built on \emph{\ru{грань}} \textquotedblleft{}an edge\textquotedblright{}: to restrict is to draw a border round a thing.
\end{cousinweb}

\begin{grammarlens}[title={What you've built}]
One more word for telling, explaining and describing at B1.
\end{grammarlens}

\subsection*{Your turn}
\begin{itemize}
\raggedright
  \item \emph{Say it:} \emph{ogranichéniye}
  \item \emph{Say it:} \emph{ogranichéniye}, once more
  \item \emph{Say it:} a sentence of your own with \emph{ogranichéniye}
  \item \emph{Recall:} say the Russian for citizenship, then the Russian for of the law, then say \emph{ogranichéniye} again
\end{itemize}

\begin{tcolorbox}[breakable,colback=teal!4,colframe=teal!35!black,title={Before you move on}]
What does \ru{ограничение} mean? (\textquotedblleft{}A restriction, a limit\textquotedblright{}.) Say it once more.
\end{tcolorbox}
